% !TEX root = ../matan.tex

\section{Достаточное условие дифференцируемости функции.}

\begin{Theorem}{Достаточное условие дифференцируемости}{}
	Пусть $\Omega\subset\R^d$ открыто, $f\colon\Omega\to\R$, $x\in\Omega$. Если все частные производные $\dfrac{\partial f}{\partial x_j}$ ($j=1,\dots,d$) существуют в некотором шаре $B(x,\eps)\subset\Omega$ и непрерывны в самой точке $x$, то $f$ дифференцируема в точке $x$, причём
	\[
	d_xf(h)=\sum_{j=1}^{d}\frac{\partial f}{\partial x_j}(x)\,h_j=\langle\nabla f(x),h\rangle.
	\]
	Функцию, у которой все частные производные существуют и непрерывны на $\Omega$, называют \emph{непрерывно дифференцируемой} ($f\in C^1(\Omega)$); по теореме такая функция дифференцируема в каждой точке $\Omega$.
\end{Theorem}

\begin{proof}
	Проведём доказательство для $d=2$; общий случай получается тем же приёмом (последовательная замена координат по одной). Пусть $x=(x_1,x_2)$, $\xi=(\xi_1,\xi_2)\in B(x,\eps)$. Шар выпукл, поэтому вместе с $x$ и $\xi$ ему принадлежит и промежуточная точка $(x_1,\xi_2)$, а также все точки отрезков, по которым ниже применяется теорема Лагранжа.

	\textbf{Шаг 1. Разбиение приращения.} Меняем координаты поочерёдно:
	\[
	f(\xi_1,\xi_2)-f(x_1,x_2)
	=\bigl[f(\xi_1,\xi_2)-f(x_1,\xi_2)\bigr]+\bigl[f(x_1,\xi_2)-f(x_1,x_2)\bigr].
	\]

	\textbf{Шаг 2. Теорема Лагранжа по каждой переменной.} В первой скобке меняется только первый аргумент. Функция одной переменной $\vphi(t)=f(t,\xi_2)$ дифференцируема на отрезке с концами $x_1$ и $\xi_1$ (там существует $\partial_1 f$, а её значение есть $\vphi'(t)=\partial_1 f(t,\xi_2)$), поэтому по теореме Лагранжа найдётся точка $s_1$ между $x_1$ и $\xi_1$ такая, что
	\[
	f(\xi_1,\xi_2)-f(x_1,\xi_2)=\frac{\partial f}{\partial x_1}(s_1,\xi_2)\,(\xi_1-x_1).
	\]
	Аналогично, $\psi(t)=f(x_1,t)$ даёт точку $s_2$ между $x_2$ и $\xi_2$ такую, что
	\[
	f(x_1,\xi_2)-f(x_1,x_2)=\frac{\partial f}{\partial x_2}(x_1,s_2)\,(\xi_2-x_2).
	\]

	\textbf{Шаг 3. Выделение линейной части.} Прибавим и вычтем значения частных производных в самой точке $x$:
	\[
	f(\xi)-f(x)=\frac{\partial f}{\partial x_1}(x)(\xi_1-x_1)+\frac{\partial f}{\partial x_2}(x)(\xi_2-x_2)+r_1+r_2,
	\]
	где
	\[
	r_1=\Bigl(\tfrac{\partial f}{\partial x_1}(s_1,\xi_2)-\tfrac{\partial f}{\partial x_1}(x)\Bigr)(\xi_1-x_1),\qquad
	r_2=\Bigl(\tfrac{\partial f}{\partial x_2}(x_1,s_2)-\tfrac{\partial f}{\partial x_2}(x)\Bigr)(\xi_2-x_2).
	\]
	Линейная по $\xi-x$ часть имеет вид $L(\xi-x)=\sum_{j=1}^{2}\partial_j f(x)(\xi_j-x_j)=\langle\nabla f(x),\xi-x\rangle$.

	\textbf{Шаг 4. Остаток есть $o(|\xi-x|)$.} При $\xi\to x$ точки $(s_1,\xi_2)$ и $(x_1,s_2)$ также стремятся к $x$ (поскольку $s_1$ лежит между $x_1$ и $\xi_1$, а $s_2$ --- между $x_2$ и $\xi_2$). По непрерывности частных производных в точке $x$
	\[
	\alpha_1:=\frac{\partial f}{\partial x_1}(s_1,\xi_2)-\frac{\partial f}{\partial x_1}(x)\xrightarrow[\xi\to x]{}0,\qquad
	\alpha_2:=\frac{\partial f}{\partial x_2}(x_1,s_2)-\frac{\partial f}{\partial x_2}(x)\xrightarrow[\xi\to x]{}0.
	\]
	Так как $|\xi_1-x_1|\le|\xi-x|$ и $|\xi_2-x_2|\le|\xi-x|$,
	\[
	|r_1+r_2|\le\bigl(|\alpha_1|+|\alpha_2|\bigr)\,|\xi-x|=o(|\xi-x|).
	\]
	Следовательно,
	\[
	f(\xi)=f(x)+L(\xi-x)+o(|\xi-x|),
	\]
	то есть $f$ дифференцируема в точке $x$, а её дифференциал равен $d_xf=L$, $d_xf(h)=\sum_{j}\partial_j f(x)h_j$.
\end{proof}

\begin{Remark}{Существования частных производных недостаточно}{}
	Само по себе существование частных производных в точке не гарантирует дифференцируемости. Так, для $f(x_1,x_2)=\sqrt{|x_1 x_2|}$ обе частные производные в нуле существуют и равны нулю (вдоль каждой координатной оси $f\equiv0$), а сама $f$ непрерывна; тем не менее $f$ \emph{не} дифференцируема в $(0,0)$. В самом деле, будь она дифференцируема, её дифференциал был бы нулевым (по необходимому условию равен матрице из нулевых частных производных), и тогда $f(\xi)=o(|\xi|)$; но вдоль прямой $x_1=x_2=t$ имеем $f(t,t)=|t|$ и
	\[
	\frac{f(t,t)}{|(t,t)|}=\frac{|t|}{\sqrt2\,|t|}=\frac{1}{\sqrt2}\not\to0.
	\]
	Причина в том, что частные производные разрывны в нуле; именно их непрерывность и требуется в теореме.
\end{Remark}

\begin{Remark}{Векторный случай}{}
	Для отображения $f=(f_1,\dots,f_m)\colon\Omega\to\R^m$ достаточное условие применяется покомпонентно: если все частные производные $\partial f_i/\partial x_j$ существуют в окрестности точки $x$ и непрерывны в $x$, то каждая $f_i$ дифференцируема в $x$, а значит (по теореме о покомпонентной дифференцируемости, вопрос~38) дифференцируемо и само $f$, причём $d_xf=J_f(x)$.
\end{Remark}
